\begin{figure}[htbp]
    \centering
    \includegraphics[width=\linewidth]{figure1_forward_displacement_and_tracking_rmse.pdf}
    \caption{
    所提 IPM 控制器与 SQP 基线在 10 次独立闭环实验中的前向位移与跟踪 RMSE 对比。误差棒表示重复实验的一倍标准差。
    该图表明，本文方法不仅获得了显著更大的前向推进，而且以更好的重复性降低了跟踪误差。具体而言，IPM 方案能够稳定地将指令前向运动转化为真实物理推进，而 SQP 基线基本停滞，并表现出系统性更差的跟踪质量。}
    \label{fig:forward_rmse}
\end{figure}

\begin{figure}[htbp]
    \centering
    \includegraphics[width=\linewidth]{figure2_guard_and_max_still_boxplots.pdf}
    \caption{
    所提方法与 SQP 基线在 10 次重复实验中的 branch-guard 触发次数以及最大停滞时长分布。箱线图给出了中位数、四分位区间、须和离群点。
    该图直接证明了本文方法在稳定性与安全性上的优势：branch-guard 触发次数下降为零，机器人持续保持运动；而基线方法则表现为高频 guard 触发和长时间停滞。这说明 IPM 与分支一致性建模的结合消除了此前必须依赖 runtime 后置拦截才能止损的失效模式。}
    \label{fig:guard_still}
\end{figure}

\begin{figure}[htbp]
    \centering
    \includegraphics[width=0.82\linewidth]{figure3_mpc_solve_time_with_realtime_budget.pdf}
    \caption{
    所提方法与 SQP 基线的平均 MPC 求解时间对比，误差棒表示 10 次重复实验的一倍标准差。红色虚线表示 10~ms 的实时预算上限。
    尽管 IPM 的计算开销略高于 SQP，但其平均求解时间约为 4.8~ms，仍然远低于实时控制红线。这表明本文方法是在不牺牲在线部署能力的前提下，获得了更好的物理一致性和更强的闭环鲁棒性。}
    \label{fig:solve_time}
\end{figure}

\subsection{运动性能与跟踪精度}

图~\ref{fig:forward_rmse} 汇总了所提方法与 SQP 基线在 10 次独立闭环实验中的行走结果。两者在定性与定量上都呈现出明显分离。本文方法获得了 0.065~m 的平均前向位移，而 SQP 基线几乎完全停滞，平均位移仅为 $-0.002$~m。与此同时，跟踪 RMSE 也从基线的 0.242~m 降低到了本文方法的 0.175~m。

这一结果的重要性在于两个层面。首先，它说明从 SQP 切换到 IPM 并非单纯的数值微调，而是恢复了具有物理意义的前向行走能力。其次，更低的跟踪误差意味着机器人不再需要用控制努力去补偿内部不一致的关节空间解，指令运动能够更直接地转化为 Isaac Sim 中的真实机体推进。

需要指出的是，方差在这里必须被谨慎解释。SQP 基线的前向位移分布较窄，并不意味着它更鲁棒，而是意味着它以较高一致性“稳定失败”。相对地，本文方法不仅实现了正向推进，而且试验间离散度较小，这才是具有实际意义的鲁棒性指标。换言之，本文方法不仅平均表现更好，而且是可重复地更好。

因此，图~\ref{fig:forward_rmse} 证明了这样一个核心结论：一旦将机构分支一致性以硬约束形式写入最优控制问题，并使用能够真实处理该约束的 IPM 求解器，六足机器人就能够恢复稳定且有效的前向行走能力。

\subsection{约束满足与系统安全性}

图~\ref{fig:guard_still} 展示了本文方法在安全性与稳定性方面的收益。最显著的现象是：本文方法在全部 10 轮实验中 branch-guard 触发次数为零，而 SQP 基线的平均触发次数达到 275.9 次/轮。同时，最大停滞时长从 5.64~s 降至 0.64~s，trunk sink 触发率也从 10\% 降为 0。

这些统计量清楚地反映出求解器层面修复的实际意义。在 SQP 基线下，控制器会反复生成必须被 runtime 安全逻辑纠正或直接拒绝的输出。这种行为的问题不只在于运动效率下降，更在于它说明优化器仍然频繁访问物理上错误的状态分支，使 runtime 层被迫承担“灾后处理”的角色，而不是轻量的安全监督者。

相比之下，本文方法在测试范围内已经不再依赖这类后置拦截机制。这说明 \emph{branch-consistency constraint} 不再只是一个理论上存在的条件，而是在求解器内部真正塑造了 walking phase 的可行域。物理上错误的 tarsus 分支在尚未形成可执行关节命令之前，就已经被排除在外。从安全性的角度看，这是从“事后拒绝危险输出”向“事前消除危险解分支”的根本转变。

停滞时长的显著下降进一步支持了这一解释。基线中的长时间静止并不是“谨慎控制”的体现，而是优化不一致与 guard 高频介入的外部表现。一旦错误分支被从求解过程中去除，机器人便能保持连续且平滑的运动过程。图~\ref{fig:guard_still} 因而不仅证明了任务表现的提升，也证明了闭环规则性与系统安全性的实质改善。

\subsection{实时性与计算代价}

图~\ref{fig:solve_time} 比较了两种方法的平均 MPC 求解时间。正如预期，本文的 IPM 方案在计算上略重于 SQP 基线，平均求解时间分别为 4.82~ms 与 3.56~ms。然而，两者都明显低于 10~ms 的实时预算上限，说明本文方法在给定硬件和仿真链路上完全满足在线控制的时间要求。

这一结果对工程部署尤为关键。一个常见担忧是：采用更严格处理约束的内点法，可能会用鲁棒性换取不可接受的延迟。本文实验并不支持这一担忧。IPM 带来的额外计算开销在绝对量级上较小，而控制器依然保持在安全的实时运行区间之内。

更重要的是，这一额外算时必须结合行为收益来理解。SQP 的确更快，但它更快地求解了一个“行为意义上错误的问题”，因为错误分支仍然留在有效搜索空间中，最后还必须依赖 runtime guard 进行补救。本文方法则用略高的算时换来了零 guard、零 reject、显著更强的前向推进以及更低的跟踪误差。从系统层面看，这是一笔非常划算的交换。

因此，图~\ref{fig:solve_time} 支持了本文的另一项核心论断：更强的物理一致性和更高的闭环可靠性，并不一定需要以牺牲实时性为代价。IPM 在这里不是“过度昂贵的学术求解器”，而是一个完全可部署的实时 MPC 后端。

\subsection{长航时扩展验证与鲁棒性}

前述受控 A/B 对照已经建立了求解器可见分支一致性的因果收益。接下来，我们进一步考察：在求解器层面的根因修复已经被证明有效之后，控制器在更长时间尺度下是否仍能维持稳定行走。需要强调的是，这里的 endurance 结果并不是为了替代短程 A/B 对照，而是为了将本文的贡献从“因果诊断正确”进一步扩展到“系统级持续有效”。

\begin{figure*}[t]
    \centering
    \includegraphics[width=0.92\textwidth]{endurance_timeseries.pdf}
    \caption{
    最终 60~s endurance 评测中记录到的闭环时序曲线。上图为前向位移，显示机器人在采集到的闭环窗口内持续推进并最终达到约 0.638~m。中图为偏航漂移，表明航向误差始终保持在工程上可控的范围内，而没有演化为不稳定长测中常见的大幅发散。下图为机身高度变化，整体波动仅在厘米量级，总极差约为 0.011~m。三条曲线共同说明：本文方法在长航时运行中同时保持了持续推进、受控偏航和稳定机身姿态。}
    \label{fig:endurance_timeseries}
\end{figure*}

正如图~\ref{fig:endurance_timeseries} 所示，长航时响应呈现出明显的“有界稳定”特征，而不是后期失控：前向位移持续积累，偏航保持受控，机身高度波动也维持在较小范围内。在 5 次重复的 20~s endurance 评测中，本文方法获得了 0.306~m 的平均 Isaac 前向位移，标准差仅为 0.012~m；平均绝对偏航漂移为 5.85$^\circ$，平均最大停滞时长为 0.66~s，平均 MPC 求解时间为 4.62~ms。更重要的是，这一批次中的全部实验都保持了零 branch-guard、零 backend reject，且没有出现 trunk sink 或 weird pose。该结果表明，本文方法不仅能在短程测试中“做对一次”，还能够在中等时长的重复实验中以较低离散度持续保持正确前向行走。

在进一步的 60~s 长航时评测中，机器人实现了 0.638~m 的前向位移，偏航漂移为 8.27$^\circ$，机身高度变化范围仅 0.011~m，而平均 MPC 求解时间仍保持在 4.63~ms。整个过程中，mode schedule 始终处于预期的行走模式，没有出现 branch-guard 触发、backend reject 或 fault reset。这些结果说明，即使将测试时长扩展到 20~s 评测的三倍，控制器依然能够将航向误差和机身高度波动约束在可控范围内。

综上，endurance 结果表明，本文提出的求解器级修复不仅在短程 acceptance setting 下成立，而且能够向更长时域推广。虽然前向位移并未随时间呈现完全线性的增长，但控制器保留了两个最关键的可部署性质：其一，求解器吞吐始终维持在实时区间；其二，机器人未再进入原始 SQP 流程中常见的错误分支切换或姿态崩塌模式。从这一意义上说，长航时实验强化了本文的实践价值：所提方法不仅数值上正确，而且系统层面具有持续运行的能力。

\subsection{讨论与更广泛的意义}

本文更深层的学术贡献，在于揭示了高维多足系统中 SQP 型 MPC 的一种隐蔽但后果严重的失效模式。对于六足机器人而言，接触冗余、混杂步态切换以及多种关节分支可能性共同构成了一个高度复杂的优化地形。在这样的地形上，任务空间相似性可能掩盖关节空间物理语义的不一致。于是，warm-start 驱动的 SQP 可能收敛到一个在线性化意义下“似乎可接受”的局部极小值，但该极小值对应的其实是机构上错误的分支。

这一失效模式尤其危险，因为它往往不会立刻以数值发散的形式暴露出来。相反，控制器可能在指标上看起来“还算正常”，却持续生成必须在 runtime 层面被抑制的命令。也就是说，控制器表面健康，实则在系统层面不断制造后置止损压力。

本文提出的修复方案从架构上改变了这一局面。通过将 \emph{branch-consistency constraint} 作为硬约束与 IPM 后端结合，我们不再把物理分支有效性视为一个优化之后再判断的“偏好”，而是将其直接提升为求解器内部的可行性条件。这个变化在概念上非常重要。runtime guard 能止损，但不能修复一个错误建模的最优控制问题；而可行域塑形则是在病灶源头消除问题。

当然，本文并不声称已经在任意地形、任意扰动和任意接触不确定性下完成了全局安全证明。但本文确实为这样的安全证明建立了更坚实的前提。一旦物理上不可能或不期望的分支已经在求解器层面被排除，后续关于安全性的理论分析就可以建立在一个与真实机构语义一致的优化域之上。从这个意义上说，本文方法为面向复杂足式机器人控制的\emph{可证明安全性}提供了一种更有原则的新路径：不是在优化后叠加越来越激进的拦截规则，而是从一开始就让危险分支在优化问题中结构性地不存在。

随着机器人面向更复杂地形和更高机动性发展，这一思路的重要性只会进一步上升。未来评价一个求解器的优劣，不能只看目标值是否下降、算时是否够快，还必须看它是否保留了机构与物理层面的正确语义。本文结果表明，当内点法与具有明确物理意义的硬约束相结合时，它为多足机器人实时 MPC 提供了一条兼具可扩展性、稳定性与安全解释能力的方向。
